\section{Условие}

{\bfseries Цель работы:}

Приобретение практических навыков управления процессами и организации
межпроцессного взаимодействия с использованием каналов.

\vspace{10pt}

{\bfseries Задание:}

Родительский процесс должен создать дочерний процесс. Родительский и
дочерний процессы должны быть представлены различными программами.

Первой строкой пользователь вводит в консоль родительского процесса имя
файла, которое передаётся дочернему процессу при его создании.

Родительский процесс передаёт команды пользователя дочернему процессу
через канал, связанный со стандартным входным потоком дочернего процесса.
Результаты работы дочерний процесс записывает в файл, имя которого было
получено от родительского процесса.

Необходимо обрабатывать ошибки используемых системных вызовов и
освобождать используемые ресурсы.

\vspace{10pt}

{\bfseries Вариант:} 2

Пользователь вводит команды, содержащие произвольное количество чисел
типа \texttt{float}, разделённых пробелами. Каждая команда заканчивается
символом перевода строки. Родительский процесс передаёт введённые данные
дочернему процессу. Дочерний процесс вычисляет сумму чисел каждой команды
и записывает результат в файл.